\chapter{Chủ Nhiệm Và Không Khí}
\chapterintro{Giờ sinh hoạt rất cần lời có duyên}{Nói đúng chất giọng ở giờ chủ nhiệm giúp lớp mở lòng hơn nhiều.}

\section{Tuần này lớp mình đáng khen ở điểm nào trước?}
\begin{manthoai}{Màn thoại}
\textbf{Thầy/Cô:} Tuần này lớp mình đáng khen ở điểm nào trước?\\
\textbf{Học sinh:} Chưa bị nhắc đồng phục nhiều ạ.\\
\textbf{Thầy/Cô:} Tốt. Có nghĩa là mình có tiến, không phải chỉ có lỗi.
\end{manthoai}
\begin{dungkhinào}
Rất hợp để mở giờ chủ nhiệm theo hướng tích cực, tránh vừa vào đã kiểm điểm.
\end{dungkhinào}
\begin{nhipthem}
Cho 2 đến 3 ý khen trước rồi mới đi vào phần cần sửa.
\end{nhipthem}
\ngancach

\section{Ai có lỗi nhỏ mà muốn sửa lớn?}
\begin{manthoai}{Màn thoại}
\textbf{Thầy/Cô:} Ai có lỗi nhỏ mà muốn sửa lớn?\\
\textbf{Học sinh:} Nghe cũng áp lực ạ.\\
\textbf{Thầy/Cô:} Không áp lực. Chỉ là mình sửa tử tế hơn tuần trước thôi.
\end{manthoai}
\begin{dungkhinào}
Dùng khi muốn nói về thay đổi hành vi mà không biến buổi sinh hoạt thành bản án.
\end{dungkhinào}
\begin{nhipthem}
Cho học sinh tự nêu một điều sẽ sửa trong tuần tới. Cách này khiến phần chủ nhiệm bớt áp đặt.
\end{nhipthem}
\ngancach

\section{Lớp mình cần bớt điều gì để vui hơn?}
\begin{manthoai}{Màn thoại}
\textbf{Thầy/Cô:} Lớp mình cần bớt điều gì để vui hơn?\\
\textbf{Học sinh:} Bớt drama ạ.\\
\textbf{Thầy/Cô:} Câu này đáng viết lên bảng đấy.
\end{manthoai}
\begin{dungkhinào}
Hợp với giờ sinh hoạt cần xả áp và nhìn lại tập thể nhẹ nhàng.
\end{dungkhinào}
\begin{nhipthem}
Sau đó có thể hỏi thêm: \textit{``Vậy mình thêm điều gì để lớp dễ chịu hơn?''}
\end{nhipthem}
\ngancach

\section{Cuối tháng này lớp muốn tự hào vì điều gì?}
\begin{manthoai}{Màn thoại}
\textbf{Thầy/Cô:} Cuối tháng này lớp muốn tự hào vì điều gì?\\
\textbf{Học sinh:} Muốn đỡ bị nhắc hơn ạ.\\
\textbf{Thầy/Cô:} Tốt. Có đích rồi thì mình đi cho ra đích đó.
\end{manthoai}
\begin{dungkhinào}
Hợp khi muốn chuyển giờ chủ nhiệm từ chỗ xử lý sự vụ sang chỗ tạo mục tiêu chung.
\end{dungkhinào}
\begin{nhipthem}
Chỉ nên lấy 1 đến 2 mục tiêu, đừng tham nhiều làm lớp khó nhớ và khó làm.
\end{nhipthem}
\ngancach

\section{Lớp mình cần thêm một thói quen nhỏ nào?}
\begin{manthoai}{Màn thoại}
\textbf{Thầy/Cô:} Lớp mình cần thêm một thói quen nhỏ nào để tuần sau dễ thở hơn?\\
\textbf{Học sinh:} Vào lớp nhanh hơn ạ.\\
\textbf{Thầy/Cô:} Một thói quen nhỏ mà cứu được nhiều phút lớn đấy.
\end{manthoai}
\begin{dungkhinào}
Dùng khi muốn nói chuyện nề nếp mà không biến buổi chủ nhiệm thành buổi phán xét.
\end{dungkhinào}
\begin{nhipthem}
Câu hỏi này hiệu quả vì nó kéo lớp về thứ có thể làm được ngay trong tuần tới.
\end{nhipthem}
